\section{Exact certificates for the scalar bounds}
\label{app:certificates}

This appendix proves Lemma~\ref{lem:scalar-bounds}.  All terminating
decimals in this appendix denote exact rational numbers.  The argument
uses two finite tables of rational polynomial minorants.  We give the
coefficient rule and the Bernstein conversion formula explicitly, so the
tables can be verified using only rational arithmetic; the accompanying
programs are supplementary implementations of these formulas.

Throughout, $X=1-x$, $0\le x\le1$, $0\le w\le X$, $0\le r\le1$, and
$s\in\{0,1\}$.  We recall the notation
\begin{align*}
 P&=\int_0^1t e^{-tx}\,dt,&
 A&=3-2\int_0^1e^{-tx}\,dt,&
 m&=A-1-2Px,\\
 C&=1-sP,& a&=A+2Pw,& b&=(2-s)P(x+w),\qquad z=Pw^2,\\
 k_s&=(1-s)\frac{1-\sqrt{1-P}}{1+\sqrt{1-P}},&
 K&=C-k_s-br+1.6z,\\
 L&=(C-m-b)r+.4z,&
 M&=(1-r)a+CPr,& G&=M-L.
\end{align*}
Also set
\[
 d_0=.52,\qquad d_1=.4,\qquad
 H_0=.394-.42X^2+.04X^3,\qquad H_1=.25,\qquad J_s=CH_s.
\]
The elementary integral bounds give $0<P\le.5$ and $1\le A\le2$;
the latter follows from $1-e^{-tx}\le tx$.  Direct integration gives
\begin{equation}
 (1+x)(A-1)=2x(1-P).
 \label{eq:cert-integral-identity}
\end{equation}

\subsection{Rational approximations and their errors}
\label{sec:cert-interpolation}

Use the following polynomials in $X$, with coefficients displayed in
ascending order:
\begin{equation}
\begin{aligned}
 p(X)={}&.26424112+.16053514X+.05751895X^2\\
        &\hspace{12mm}+.01313705X^3+.00456774X^4,\\
 v(X)={}&1.73575888-.52832269X-.16191093X^2\\
        &\hspace{12mm}-.03442938X^3-.01109588X^4,\\
 k(X)={}&.0767+.0564X+.0142X^2+.0246X^3.
\end{aligned}
\label{eq:cert-interpolants}
\end{equation}
We will prove
\begin{equation}
 |p-P|\le\epsilon_P:=.0000075,\qquad
 |v-A|\le2\epsilon_P=.000015,\qquad k_0\le k.
 \label{eq:cert-approximation-errors}
\end{equation}

For the first two claims, the integral Taylor series are
\[
 P(x)=\sum_{n\ge0}\frac{(-x)^n}{n!(n+2)},\qquad
 A(x)=3-2\sum_{n\ge0}\frac{(-x)^n}{n!(n+1)}.
\]
On $0\le x\le1$ the magnitudes of the summands decrease.  Thus the
partial sums through orders $12$ and $13$ enclose each defining
integral.  Evaluating these rational enclosures at
$X=0,1/4,1/2,3/4,1$ proves that the nodal errors of the displayed
$p,v$ are, respectively, less than
$2.343560\cdot10^{-9}$ and $2.344324\cdot10^{-9}$.  In particular,
each is less than $4\cdot10^{-9}$.  This nodal calculation can be
performed without evaluating any transcendental number: at each node
substitute $x=1-X$ into the two finite sums just given and compare
both rational endpoints with the displayed polynomial value.

Let $I_P,I_A$ be the exact degree-four interpolants at these five nodes.
The sum of the absolute values of the Lagrange basis polynomials is at
most
\[
 \sum_{i=0}^4\frac{4^4}{i!(4-i)!}=\frac{512}{3},
\]
because every numerator factor has absolute value at most one.
Differentiating with respect to $X$ gives the fifth-derivative bounds
$1/7$ for $P$ and $1/3$ for $A$.  The absolute node product, with
$u=|X-1/2|$, is
\[
 u(1/4-u^2)|1/16-u^2|.
\]
For $u\le1/4$ it is at most $1/256$.  For $1/4\le u\le1/2$, the
product of the two quadratic factors is at most
$((1/4-1/16)/2)^2=9/1024$, so the entire expression is at most
$9/2048$.  Both bounds are less than $.0045$.  The interpolation
remainder formula therefore gives
\begin{align*}
 |p-P|&\le\frac{512}{3}(4\cdot10^{-9})+\frac{.0045}{840}
       =\frac{31709}{5250000000}<\epsilon_P,\\
 |v-A|&\le\frac{512}{3}(4\cdot10^{-9})+\frac{.0045}{360}
       =\frac{9887}{750000000}<2\epsilon_P.
\end{align*}

The coefficient signs and endpoint values give
\begin{equation}
\begin{gathered}
 .26424112\le p\le.5,\qquad 1\le v\le1.73575888,\\
 .0767\le k\le.1719,\qquad .014\le H_0\le.394.
\end{gathered}
\label{eq:cert-polynomial-ranges}
\end{equation}
For $H_0$, use $H_0'(X)=X(-.84+.12X)\le0$.
The first row of Table~\ref{tab:cert-first} below proves
\begin{equation}
 T_k:=4k-(p+\epsilon_P)(1+k)^2>0.
 \label{eq:cert-k-envelope}
\end{equation}
Since $0<k<1$ and $t\mapsto4t/(1+t)^2$ is increasing on $(0,1)$,
while $4k_0/(1+k_0)^2=P$, inequality~\eqref{eq:cert-k-envelope}
and $P\le p+\epsilon_P$ imply $k_0\le k$.  This proves all of
\eqref{eq:cert-approximation-errors}.

Define polynomial versions of the quantities used in the following reductions by
\begin{equation}
\begin{gathered}
 c=1-sp,\qquad N=c-(1-s)k,\qquad l=c+1-v+spx,\\
 d=d_s,\qquad h=1.6+.4d,\qquad
 u_0=d/h,\qquad u_1=(2-s)(1+d)/(2h).
\end{gathered}
\label{eq:cert-polynomial-notation}
\end{equation}
The polynomial families in the reductions are
\begin{align*}
 T_1&=1-p-1.12p(1+x)+.4p^2X(1+x),\\
 T_2&=((1-p)^2-.38pX+px)(1+x)-2x(1-p),\\
 Q_0(W)&=N-dv+pW(hW-2d),\\
 Q_1&=N-(2-s)px-d(cp-l)-hp u_1^2,\\
 Y&=N-p\bigl((2-s)x+(2-s)^2/6.4\bigr),\\
 Y'&=N+p(1.6X^2-2)\qquad(s=0).
\end{align*}
The notation $T(P,A)$ means replacement of $p,v$ by $P,A$, keeping
$X,k$ and the rational constants fixed.


\subsection{Reduction of the first, second, and final scalar bounds}

We first establish the first scalar inequality
\begin{equation}
 m+b-.4z/r\le\frac{C(1-P)}{1-Pr}\qquad(0<r\le1).
 \label{eq:cert-first-scalar}
\end{equation}
Put $t=1/r-P>0$.  By~\eqref{eq:cert-integral-identity}, the right
side of~\eqref{eq:cert-first-scalar} minus its left side equals
\begin{align*}
 &C(1-P)-\frac{2x(1-P)}{1+x}+sPx-(2-s)Pw+.4P^2w^2\\
 &\hspace{35mm}+P\left(\frac{C(1-P)}t+.4w^2t\right).
\end{align*}
Applying AM--GM to the last parentheses gives the lower bound
\begin{equation}
 C(1-P)-\frac{2x(1-P)}{1+x}+sPx
 -(2-s-2\sqrt{.4C(1-P)})Pw+.4P^2w^2.
 \label{eq:cert-amgm}
\end{equation}
For $s=0$ the coefficient $2-2\sqrt{.4(1-P)}$ is at most $1.12$.
The resulting quadratic decreases for $0\le w\le X$: its derivative
is at most $-1.12P+.8P^2<0$.  Substituting $w=X$ into this lower
bound gives $X T_1(P)/(1+x)$, where $T_1$ is defined above.  For
$s=1$, the coefficient $1-2\sqrt{.4}(1-P)$ is at most $.38$.
Dropping the nonnegative quadratic term and using $w\le X$ gives
$T_2(P)/(1+x)$.  The numerical coefficient comparisons also follow
by squaring positive quantities: $.2>.44^2$ and
$.4>.62^2$.  Table~\ref{tab:cert-first}, with the substitution
budget proved below, gives $T_1(P),T_2(P)>0$.
At $X=0$, the extra factor $X$ in the $s=0$ reduction is zero,
which still proves the required non-strict inequality.

For $r>0$, \eqref{eq:cert-first-scalar} implies
\[
 L\ge Cr\left(1-\frac{1-P}{1-Pr}\right)\ge0.
\]
For $r=0$, one has $L=.4z\ge0$ directly.

To check $K\ge d_sG$, write $d=d_s$ and $h=1.6+.4d$.
The difference $K-dG$ is affine in $r$.  With
$N_P=C-(1-s)k$, replacing $k_s$ by $(1-s)k$ only decreases this
difference.  Put $\underline K=N_P-br+1.6z\le K$.  The endpoint
expressions after this replacement are
\begin{align}
 (\underline K-dG)|_{r=0}&=N_P-dA+Pw(hw-2d),
 \label{eq:cert-r0}\\
 (\underline K-dG)|_{r=1}&=Q_1(P,A)+hP(w-u_1)^2,
 \label{eq:cert-r1}
\end{align}
where
\[
 u_0=d/h,\qquad u_1=\frac{(2-s)(1+d)}{2h}
\]
and $Q_1$ is defined above.  For $s=1$ we may use the unrestricted
minimum $w=u_0$ in~\eqref{eq:cert-r0}.  For $s=0$, the constrained
minimum is at $w=X$ if $X\le u_0$, and at $w=u_0$ otherwise.
Equation~\eqref{eq:cert-r1} is bounded below by $Q_1(P,A)$ for
both values of $s$.  These are precisely the $Q_0,Q_1$ checks in
Table~\ref{tab:cert-first}.

We next reduce the final two scalar bounds.  Recall
\[
 \kappa_s(q)=C\frac{1-q}{1-Pq}-(1-q),\qquad 0\le q\le1.
\]
For $s=1$, $k_s=0$ and $\kappa_s(q)\le0$.  Thus both relevant
left sides are bounded below by $C+1.6z-b$.  Its unrestricted
minimum in $w$ is the true-parameter version of $Y$ above, so the
rows $Y-H_0$ and $Y-.25$ suffice.

For $s=0$ and $X\ge.5$, use $q\le1$ and $\kappa_0(q)\le k_0$.
Again a common lower bound is $1-k_0+1.6z-b$, whose unrestricted
minimum is the corresponding $Y$.  For $s=0$ and $X\le.5$, we
have $x\ge.5$ and
\[
 P(x)\le P(.5)<.47.
\]
For example the inequality $e^{-u}\le1-u+u^2/2$ yields
$P(.5)\le1/2-1/6+1/32=35/96<.47$.
It follows that $b\ge .52P/(1-P)$, because $x+w\ge.5$ and
$1-P>.53$.  Since
\[
 \kappa_0(q)=\frac{Pq(1-q)}{1-Pq}
       \le\frac{P}{1-P}(1-q),
\]
we obtain $bq+.52\kappa_0(q)\le b$.  The fourth scalar left side
is consequently at least
\[
 1-.48k_0+1.6z-b.
\]
The quadratic $1.6Pw^2-2P(x+w)$ decreases on $0\le w\le X\le.5$,
so its minimum is $P(1.6X^2-2)$.  Thus the fourth and fifth
inequalities follow, respectively, from $Y'+.52k-H_0\ge0$ and
$Y'-.25\ge0$ after substitution.  In the former expression the
net coefficient of $k$ is $-.48$, so replacing $k_0$ by the upper
bound $k$ is conservative.

The last scalar bound also gives positivity of $K$: since $b\ge0$,
\[
 K=C-k_s-br+1.6z\ge C-k_s-b+1.6z\ge.25>0.
\]

\subsection{The first coefficient table and its exact verification}

Replacing $p,v$ by $P,A$ in these expressions changes their values
by less than
\begin{equation}
 \sigma=.0002.
 \label{eq:cert-sigma}
\end{equation}
For completeness, the following derivative bounds hold on the line
segment between $(p,v)$ and $(P,A)$, using $0\le p,P\le.5$,
$1\le v,A\le2$, and $0\le W\le1$:
\[
\begin{array}{c|rr|r}
 T&|\partial_pT|&|\partial_vT|&
 10^6\,|T(p,v)-T(P,A)|\text{ upper bound}\\ \hline
 T_1&3.64&0&27.3\\
 T_2&8.76&0&65.7\\
 Q_0(W)&3.848&.52&36.66\\
 Q_1&4.848&.52&44.16\\
 Y&2.625&0&19.6875\\
 Y'&2&0&15
\end{array}
\]
Here $T_{1,p}=-1-1.12(1+x)+.8pX(1+x)$, with
$X(1+x)\le1$, and
$T_{2,p}=[-2(1-p)-.38X+x](1+x)+2x$.
For $Q_0$, use $h\le1.808$ and $2d\le1.04$.
For $Q_1$, $(cp-l)_p=1-2sp+sX$ has absolute value at most $2$,
$(cp-l)_v=1$, and $u_1\le1$.
The last two derivative bounds follow directly from their displayed
linear dependence on $p$.  These estimates prove the claimed budget
even on the full interval, rather than only the listed subintervals.
Adding $.52k$ or subtracting $H_0,.25$ does not change this budget.
The row for $T_k$ already includes $\epsilon_P$ explicitly and does
not need this additional substitution estimate.

Here is the finite verification mechanism used in both coefficient
tables.  A row with scale $10^a$ and integers
$\lambda_0,\ldots,\lambda_n$ asserts
\begin{equation}
 T(X)\ge t(X):=10^{-a}\sum_{j=0}^n\lambda_jX^j
       \qquad(0\le X\le1).
 \label{eq:cert-minorant}
\end{equation}
To verify it, expand $T=\sum_jc_jX^j$ and check
\begin{equation}
 10^ac_j\ge\lambda_j\quad(0\le j<n),\qquad
 10^a\sum_{j\ge n}\min(0,c_j)\ge\lambda_n.
 \label{eq:cert-coefficient-rule}
\end{equation}
Positive tail terms may be discarded; for every negative tail term,
$c_jX^j\ge c_jX^n$.  This proves~\eqref{eq:cert-minorant}.

All short minorants have degree at most six.  On any listed interval
$[\alpha,\beta]$, their degree-six Bernstein coefficients are
\begin{equation}
 \mathfrak b_m(t;\alpha,\beta)=10^{-a}
 \sum_{h=0}^m\frac{\binom mh}{\binom6h}(\beta-\alpha)^h
       \sum_{j=h}^n\lambda_j\binom jh\alpha^{j-h},
       \quad 0\le m\le6.
 \label{eq:cert-bernstein-rule}
\end{equation}
Indeed, expand at $X=\alpha+(\beta-\alpha)u$ and convert ordinary
powers of $u$ to degree-six Bernstein form.  Thus if every coefficient
in~\eqref{eq:cert-bernstein-rule} is at least $b_*$, then
$t(X)\ge b_*$ throughout the interval, since the Bernstein basis is
nonnegative and sums to one.  An endpoints entry with more than two
numbers prescribes each interval between consecutive numbers.

\begingroup
\small
\setlength{\tabcolsep}{3pt}
\renewcommand{\arraystretch}{1.18}
\begin{longtable}{@{}lclc>{\raggedright\arraybackslash}p{.39\textwidth}@{}}
\caption{First coefficient table.  The scale is $10^a$, and the
 listed vector is $(\lambda_0,\ldots,\lambda_n)$.}
\label{tab:cert-first}\\
\toprule
 $T$&$s$&Endpoints and splits&$a$&$\lambda$\\
\midrule
\endfirsthead
\toprule
 $T$&$s$&Endpoints and splits&$a$&$\lambda$\\
\midrule
\endhead
\bottomrule
\endfoot
 $T_k$&--&$0,.5,1$&4&$(4,73,-384,563,-184,-67)$\\
 $T_1$&--&$0,1$&2&$(14,-17,3,3,-1)$\\
 $T_2$&--&$0,1$&2&$(13,10,-22)$\\
 $Q_0(u_0)$&1&$0,1$&2&$(1,3,-1)$\\
 $Q_1$&1&$0,1$&2&$(42,-3)$\\
 $Q_1$&0&$0,1$&2&$(5,-1)$\\
 $Q_0(u_0)$&0&$u_0,1$&2&$(-2,19,-1)$\\
 $Q_0(X)$&0&$0,u_0$&2&$(2,-6)$\\
 $Y-H_0$&1&$0,1$&2&$(3,-9,45,-2)$\\
 $Y-.25$&1&$0,1$&2&$(17,-9)$\\
 $Y-H_0$&0&$.5,1$&3&$(-165,50,575,-1)$\\
 $Y-.25$&0&$.5,1$&2&$(-3,5,15,-1)$\\
 $Y'+.52k-H_0$&0&$0,.125,.25,.5$&4&$(406,-3483,7208,1786,-3)$\\
 $Y'-.25$&0&$0,.5$&2&$(14,-38,29,-1)$\\
\end{longtable}
\endgroup

Substitution of~\eqref{eq:cert-interpolants} into the definitions
verifies~\eqref{eq:cert-coefficient-rule} for all fourteen rows of
Table~\ref{tab:cert-first}.  As additional explicit sign certificates,
the minima of the finite sets in~\eqref{eq:cert-bernstein-rule}, over
all prescribed subintervals for each row, are at least $10^{-6}$ times
the following integers, in row order:
\begin{equation}
\begin{gathered}
 114, 20000, 10000, 10000, 390000, 40000, 33818,\\
 2743, 15000, 80000, 3625, 31250, 356, 21250.
\end{gathered}
\label{eq:cert-first-bernstein-lower}
\end{equation}
Both~\eqref{eq:cert-coefficient-rule} and these lower bounds are finite
rational comparisons specified completely by the table and
\eqref{eq:cert-bernstein-rule}.  For example, the first row has
smallest Bernstein coefficient $73/640000$, and the thirteenth row
has smallest coefficient $27341/76800000$.
Consequently $T_k>0$, and every other short minorant exceeds
$\sigma=.0002$ on its indicated intervals.  After the substitution
budget, this completes the first and second groups of scalar bounds
and the final two scalar bounds.

\subsection{The third scalar bound: perturbation estimate}

It remains to prove
\begin{equation}
 ((J_s-b)r+1.6z)M+(C-k_s-J_sr)L\ge0.
 \label{eq:cert-third-goal}
\end{equation}
Since $L\ge0$ is already established, replacing $k_s$ with $(1-s)k$
only decreases the left side.  We next replace $P,A$ by $p,v$.
This replacement changes the expression by at most
\begin{equation}
 \sigma(r+w^2),\qquad \sigma=.0002.
 \label{eq:cert-third-error}
\end{equation}
Here are full bounds for this assertion.  For either endpoint choice
$(\pi,\nu)=(P,A)$ or $(p,v)$, put $c_\pi=1-s\pi$, $H=H_s$, and
write the two products as $F_1F_2+F_3F_4$, where
\begin{align*}
 F_1&=[c_\pi H-(2-s)\pi(x+w)]r+1.6\pi w^2,\\
 F_2&=(1-r)(\nu+2\pi w)+c_\pi\pi r,\\
 F_3&=c_\pi(1-Hr)-(1-s)k,\\
 F_4&=[2-\nu-s\pi X-(2-s)\pi w]r+.4\pi w^2.
\end{align*}
Writing $\rho=r+w^2$, the ranges already established imply
\begin{equation}
 |F_1|\le1.1\rho,\qquad |F_2|\le3,\qquad
 |F_3|\le1,\qquad |F_4|\le2\rho.
 \label{eq:cert-factor-bounds}
\end{equation}
For $F_1$, both nonnegative terms in its coefficient of $r$ are at
most one, while $1.6\pi\le.8$.  The expression $F_2$ is a convex
combination of two quantities between zero and three.  For $F_3$,
the lower bounds are $1-.394-.1719=.4341$ when $s=0$ and
$.5(1-.25)=.375$ when $s=1$, and the upper bound is one.
Finally the coefficient of $r$ in $F_4$ belongs to $[-1,1]$:
$0\le2-\nu\le1$ and $0\le s\pi X+(2-s)\pi w\le1$.

Let $\delta F_j$ denote the change between the two choices.  With
$\epsilon_P=.0000075$, direct subtraction gives
\begin{align*}
 |\delta F_1|&\le2\epsilon_P\rho,&
 |\delta F_2|&\le4\epsilon_P,\\
 |\delta F_3|&\le\epsilon_P,&
 |\delta F_4|&\le4\epsilon_P\rho.
\end{align*}
For the first estimate, the coefficient multiplying
$\epsilon_Pr$ is $sH+(2-s)(x+w)$, at most $2$ for $s=0$ and at
most $1.25$ for $s=1$; the remaining coefficient is $1.6$.
For the second, use the convex combination in $r$, together with
$|\delta(\nu+2\pi w)|\le4\epsilon_P$ and
\[
 |\delta(c_\pi\pi)|
   =|P-p|\,|1-s(P+p)|\le\epsilon_P.
\]
The third estimate is immediate from $0\le1-Hr\le1$.
For the fourth, use
$2+sX+(2-s)w\le4$ and $.4\le4$.
Combining these differences with~\eqref{eq:cert-factor-bounds} and
the identity $ab-a'b'=(a-a')b+a'(b-b')$ yields
\begin{align*}
 |\delta(F_1F_2+F_3F_4)|
 &\le(3\cdot15+1.1\cdot30+2\cdot7.5+30)10^{-6}\rho\\
 &=.000123\rho<\sigma\rho\qquad(\rho>0).
\end{align*}
If $\rho=0$, both expressions vanish.  Thus
\eqref{eq:cert-third-error} holds throughout the domain.  The factor
bounds hold for both choices, so this telescoping estimate already
accounts for products of perturbations.

\subsection{Bernstein expansion and the second coefficient table}

Use the polynomial quantities in~\eqref{eq:cert-polynomial-notation},
and set
\[
 j=cH_s,\qquad b'_0=(2-s)px,\qquad b'_1=(2-s)pX,
 \qquad f=2pX,\qquad Z'=pX^2.
\]
In the polynomial version of the left side of
\eqref{eq:cert-third-goal}, subtract $\sigma(r+w^2)$ and put
$w=X\eta$, $0\le\eta\le1$.  Its degree-three Bernstein
coefficients in $\eta$, for $i=0,1,2,3$, are
\begin{equation}
 B_0+B_1r+B_2r^2,
 \label{eq:cert-eta-coefficient}
\end{equation}
where, writing $e=i/3$, $e_2=i(i-1)/6$, and
$t=\mathbf{1}_{\{i=3\}}$, the exact formulas are
\begin{equation}
\begin{aligned}
 B_0={}&Z'\bigl((1.6v+.4N)e_2+1.6ft\bigr)-\sigma X^2e_2,\\
 B_1={}&(j-b'_0)v+Nl\\
       &+\bigl((j-b'_0)f-b'_1v-Nb'_1\bigr)e-b'_1fe_2\\
       &-Z'\bigl((1.6(v-cp)+.4j)e_2+1.6ft\bigr)-\sigma,\\
 B_2={}&(b'_0-j)(v-cp)-jl\\
       &+\bigl((b'_0-j)f+b'_1(v-cp)+jb'_1\bigr)e+b'_1fe_2.
\end{aligned}
\label{eq:cert-bcoeff}
\end{equation}
These formulas give all eight expansions, one for each pair $(s,i)$.
They can also be derived directly by expanding
\begin{align*}
 &\bigl([j-(2-s)p(x+w)]r+1.6pw^2\bigr)
       \bigl((1-r)(v+2pw)+cpr\bigr)\\
 &\quad +(N-jr)\bigl([l-(2-s)pw]r+.4pw^2\bigr)
       -\sigma(r+w^2)
\end{align*}
in powers of $\eta$ after $w=X\eta$, and using the degree-three
conversion rule: if $a_h$ is the coefficient of $\eta^h$, the
$i$th Bernstein coefficient is
$\sum_{h=0}^i a_h\binom ih/\binom3h$.

In all eight cases $B_0\ge0$.  For $i=0,1$, it is zero.  Otherwise
\[
 B_0=X^2\bigl(p(1.6v+.4N)-\sigma\bigr)e_2+1.6pX^2ft.
\]
Here $N\ge.5$, $v\ge1$, and $p\ge.26424112$, so the bracket is
at least $.475434016>0$ and the other term is nonnegative.

Define
\[
 D'=4B_0B_2-B_1^2,\qquad S'=B_0+B_1+B_2.
\]
For $i\le1$, positivity of $B_1$ and $B_1+B_2$ proves
nonnegativity of~\eqref{eq:cert-eta-coefficient} on $0\le r\le1$.
For $i\ge2$, any of the following is sufficient:
\begin{enumerate}
 \item $B_1,B_2\ge0$;
 \item $D'>0$;
 \item $S'>0$ and $B_0-B_2>0$.
\end{enumerate}
For the second criterion, $B_0\ge0$ and $D'>0$ force
$B_0,B_2>0$; completing the square proves positivity for every
real $r$.  For the third, use
\[
 B_0+B_1r+B_2r^2
   =rS'+(1-r)(B_0-rB_2),
\]
where $B_0-rB_2=(1-r)B_0+r(B_0-B_2)\ge0$.

Table~\ref{tab:cert-third} supplies these signs.  It uses exactly
the coefficient rule~\eqref{eq:cert-coefficient-rule} and the
degree-six Bernstein formula~\eqref{eq:cert-bernstein-rule}.

\begingroup
\small
\setlength{\tabcolsep}{3pt}
\renewcommand{\arraystretch}{1.18}
\begin{longtable}{@{}ccllc>{\raggedright\arraybackslash}p{.38\textwidth}@{}}
\caption{Second coefficient table.  A row with $i=0,1$ supplies
 the indicated certificate separately for both values of $i$.}
\label{tab:cert-third}\\
\toprule
 $s$&$i$&Endpoints and splits&$T$&$a$&$\lambda$\\
\midrule
\endfirsthead
\toprule
 $s$&$i$&Endpoints and splits&$T$&$a$&$\lambda$\\
\midrule
\endhead
\bottomrule
\endfoot
 0&0&$0,1$&$B_1$&2&$(1,90,-39)$\\
 0&1&$0,1$&$B_1$&2&$(1,41,-51,28,-1)$\\
 0&$0,1$&$0,1$&$B_1+B_2$&2&$(10,-1)$\\
 0&2&$0,.05$&$B_1$&2&$(1,-9,-98)$\\
 0&2&$0,.05$&$B_2$&2&$(9,-23)$\\
 0&2&$.05,1$&$D'$&4&$\begin{gathered}(-2,16,1176,-153,\\-1500,2590,-1000)\end{gathered}$\\
 0&3&$0,.015$&$B_1$&2&$(1,-58,-262)$\\
 0&3&$0,.015$&$B_2$&2&$(9,-5)$\\
 0&3&$.015,.25,.5,1$&$D'$&5&$\begin{gathered}(-12,1180,2070,-32420,\\44210,33740,-14000)\end{gathered}$\\
 1&0&$0,1$&$B_1$&2&$(5,30,-7)$\\
 1&0&$0,1$&$B_1+B_2$&2&$(12,-4)$\\
 1&1&$0,1$&$B_1$&2&$(5,-1)$\\
 1&1&$0,1$&$B_1+B_2$&2&$(12,-8)$\\
 1&2&$0,.2$&$B_1$&2&$(5,-16,-41)$\\
 1&2&$0,.2$&$B_2$&2&$(7,-2)$\\
 1&2&$.2,1$&$D'$&3&$(-4,17,92,87,-135)$\\
 1&3&$0,.1$&$B_1$&3&$(54,-391,-941,-900)$\\
 1&3&$0,.1$&$B_2$&2&$(7,-1)$\\
 1&3&$.1,.75$&$D'$&3&$(-4,42,194,478,-995)$\\
 1&3&$.75,1$&$S'$&2&$(12,-9)$\\
 1&3&$.75,1$&$B_0-B_2$&2&$(-8,-32,67,-1)$\\
\end{longtable}
\endgroup

In the printed row order, the Bernstein coefficients of these short
minorants, over all of each row's subintervals, are at least
$10^{-6}$ times the following integers:
\begin{equation}
\begin{gathered}
 10000, 10000, 90000, 3050, 78500, 171, 710,\\
 89250, 60, 50000, 80000, 40000, 40000, 1600,\\
 66000, 3560, 4590, 69000, 2518, 30000, 52656.
\end{gathered}
\label{eq:cert-third-bernstein-lower}
\end{equation}
The bound $90000$ in the third position applies to both indices in
that row.  These are all positive.  The smallest exact Bernstein
coefficient among the table's checks is
$193875074979/3200000000000000$, in the ninth printed row.
Thus all the required strict signs follow from the finite rational
formula~\eqref{eq:cert-bernstein-rule}.

For clarity, the negative-tail calculations needed for the four
discriminant rows are detailed next.  Let
$b_{h,j}=[X^j]B_h$ and write
\[
 c_j=[X^j]D'
     =\sum_u\bigl(4b_{0,u}b_{2,j-u}-b_{1,u}b_{1,j-u}\bigr),
\]
where coefficients outside the degree range are zero.  For a short
minorant of degree $n$, group $\min(0,c_j)$ in consecutive blocks
of five starting at $j=n$, truncating at degree $28$.  The five block
sums are bounded below by $.0001$ times the following integers:
\begin{equation}
\begin{array}{cc|c|c|rrrrr}
 s&i&n&\deg D'&1&2&3&4&5\\ \hline
 0&2&6&20&-930&-11&-1&0&0\\
 0&3&6&22&-1070&-180&-2&-1&0\\
 1&2&4&28&-1312&-18&-1&-1&-1\\
 1&3&4&28&-8240&-1560&-29&-1&-1
\end{array}
\label{eq:cert-tail-blocks}
\end{equation}
These twenty rational comparisons follow by substituting
\eqref{eq:cert-bcoeff} and~\eqref{eq:cert-interpolants} into the
displayed coefficient formula.  Their row sums are, respectively,
$-.0942,-.1253,-.1333,-.9831$, which are at least the final scaled
minorant coefficients $-.1,-.14,-.135,-.995$.
For all entries before the last, one checks simply
$10^a[X^j]T\ge\lambda_j$ as in
\eqref{eq:cert-coefficient-rule}.  The same coefficient rule verifies
the non-discriminant rows directly.

The tail cutoff is exact: the four discriminant degrees are
$20,22,28,28$, as recorded in~\eqref{eq:cert-tail-blocks}, so no
term above degree $28$ has been omitted.  More generally, from
\eqref{eq:cert-bcoeff}, the respective triples of degrees
$(\deg B_0,\deg B_1,\deg B_2)$ are
\[
\begin{array}{c|rrrr}
 &i=0&i=1&i=2&i=3\\ \hline
 s=0&(-\infty,9,9)&(-\infty,10,10)&(10,10,10)&(11,11,9)\\
 s=1&(-\infty,9,13)&(-\infty,10,13)&(10,14,13)&(11,14,12).
\end{array}
\]
Here $\deg0=-\infty$.  These bounds also justify a uniform cutoff of
$28$ for every polynomial in the two coefficient tables.

For $(s,i)=(0,2)$, the first criterion covers $X\in[0,.05]$ and
the discriminant criterion covers $X\in[.05,1]$.
For $(s,i)=(0,3)$, the corresponding split is $.015$, with the
listed additional subdivisions used only for Bernstein positivity.
For $(s,i)=(1,2)$, the split is $.2$.
For $(s,i)=(1,3)$, the first criterion covers $[0,.1]$, the
discriminant criterion covers $[.1,.75]$, and the final criterion
covers $[.75,1]$.  For $i=0,1$, the full interval is covered by
$B_1$ and $B_1+B_2$.  All endpoints belong to the stated intervals.
Thus every coefficient~\eqref{eq:cert-eta-coefficient} is nonnegative
for $0\le X,r\le1$.  Nonnegativity and unit sum of the degree-three
Bernstein basis in $\eta$ prove the desired polynomial inequality
for $0\le w\le X$; when $X=0$ one has $w=0$ directly.
Restoring the error budget~\eqref{eq:cert-third-error} proves
\eqref{eq:cert-third-goal}.  Together with the preceding reductions,
this completes the proof of Lemma~\ref{lem:scalar-bounds}.
